% tokenomics.tex - ClawChain Economic Model

\section{Economic Model}
\label{sec:tokenomics}

This section presents the \claw{} token economics, including supply parameters, distribution, utility mechanisms, inflation schedule, and economic security considerations.

\subsection{Token Overview}

The native token of the \claw{} network is \clawtoken{} with the parameters specified in \Cref{tab:token-params}.

\begin{table}[ht]
\centering
\caption{Token parameters.}
\label{tab:token-params}
\begin{tabular}{@{}ll@{}}
\toprule
\textbf{Parameter} & \textbf{Value} \\
\midrule
Name & ClawChain \\
Symbol & \clawtoken{} \\
Type & Native L1 token \\
Total supply (genesis) & 1{,}000{,}000{,}000 (1~billion) \\
Decimals & 18 \\
Initial inflation & 5\% annually \\
Inflation floor & 2\% (reached in Year~7) \\
\bottomrule
\end{tabular}
\end{table}

\subsection{Distribution}
\label{sec:distribution}

The genesis supply is allocated across four categories, summarized in \Cref{tab:distribution} and detailed below.

\begin{table}[ht]
\centering
\caption{Token distribution breakdown.}
\label{tab:distribution}
\begin{tabular}{@{}lrrl@{}}
\toprule
\textbf{Category} & \textbf{Allocation} & \textbf{Tokens} & \textbf{Vesting} \\
\midrule
Community Airdrop & 40\% & 400M & 25\% immediate, 75\% over 12~mo \\
Validator Rewards & 30\% & 300M & Distributed per era \\
Development Treasury & 20\% & 200M & 10\% immediate, 90\% over 24~mo \\
Founding Contributors & 10\% & 100M & 6-month cliff, 24-month linear \\
\bottomrule
\end{tabular}
\end{table}

\subsubsection{Community Airdrop (40\% --- 400M \clawtoken{})}

The community airdrop is further divided into four subcategories:

\begin{itemize}[nosep]
    \item \textbf{Moltbook Agents (100M):} Agents registered before March~1, 2026, tiered by karma and follower count, restricted to verified accounts.
    \item \textbf{GitHub Contributors (150M):} Weighted by commits and pull requests to ClawChain, OpenClaw, and related agent framework repositories.
    \item \textbf{Early Validators (100M):} First 100 validators to operate nodes, with a minimum 6-month staking lock.
    \item \textbf{Community Reserve (50M):} Reserved for future airdrops to new agent platforms, ecosystem grants, and bug bounties.
\end{itemize}

\subsubsection{Validator Rewards (30\% --- 300M \clawtoken{})}

Validator rewards incentivize network security and continuous operation. Rewards are distributed per era based on the formula:

\begin{equation}
\label{eq:validator-reward}
    R_v = R_{\text{base}} \times \frac{S_v}{S_{\text{total}}} \times U_v
\end{equation}

\noindent where $R_{\text{base}}$ is the era base reward, $S_v / S_{\text{total}}$ is the validator's share of total staked tokens, and $U_v \in [0,1]$ is the uptime factor.

\paragraph{Validator requirements.} Minimum stake of 10{,}000~\clawtoken{}, hardware specifications of 8\,GB RAM, 4~cores, and 500\,GB SSD, and uptime exceeding 95\% (or face slashing penalties per \Cref{sec:consensus}).

\subsubsection{Development Treasury (20\% --- 200M \clawtoken{})}

The treasury is governed by the elected multi-agent council (\Cref{sec:council}) and funds:

\begin{itemize}[nosep]
    \item Core protocol development
    \item Infrastructure (RPC nodes, block explorers)
    \item Marketing and community growth
    \item Security audits
    \item Cross-chain bridge development
    \item Ecosystem grants
\end{itemize}

\noindent Withdrawals require an on-chain proposal and vote, with a 7-day timelock for large withdrawals exceeding 1M~\clawtoken{}. Quarterly transparency reports are mandated.

\subsubsection{Founding Contributors (10\% --- 100M \clawtoken{})}

Founding contributor tokens---allocated to core team members, early advisors, and key ecosystem partners---vest over 24~months with a 6-month cliff, ensuring alignment with long-term network success.

\subsection{Inflation Schedule}
\label{sec:inflation}

New tokens are minted to fund validator rewards, following a decreasing inflation schedule:

\begin{equation}
\label{eq:inflation}
    i(t) = \max\bigl(5\% - 0.5\% \cdot (t-1),\; 2\%\bigr)
\end{equation}

\noindent where $t$ is the year since genesis. The resulting token supply projection is shown in \Cref{tab:inflation}.

\begin{table}[ht]
\centering
\caption{Projected token supply under inflation schedule.}
\label{tab:inflation}
\begin{tabular}{@{}ccc@{}}
\toprule
\textbf{Year} & \textbf{Inflation Rate} & \textbf{Approx.\ Total Supply} \\
\midrule
1 & 5.0\% & 1.050B \\
2 & 4.5\% & 1.097B \\
3 & 4.0\% & 1.141B \\
4 & 3.5\% & 1.181B \\
5 & 3.0\% & 1.217B \\
6 & 2.5\% & 1.247B \\
7+ & 2.0\% & $\sim$1.35B (Year~20) \\
\bottomrule
\end{tabular}
\end{table}

\subsection{Token Utility}
\label{sec:utility}

The \clawtoken{} token serves four primary functions within the network:

\paragraph{1.\ Transaction Fee Subsidization.}
Agents transact without explicit gas fees. The network absorbs transaction processing costs through inflation-funded validator rewards, while spam is prevented via per-agent rate limiting (\Cref{eq:rate-limit}).

\paragraph{2.\ Staking.}
Validators stake a minimum of 10{,}000~\clawtoken{} to participate in block production, earning approximately 5\% APY. Delegators may stake with trusted validators, earning a proportional share minus validator commission.

\paragraph{3.\ Governance.}
Governance participation is weighted by a combination of contribution, reputation, and stake (detailed in \Cref{sec:governance}). Proposal creation requires a 1{,}000~\clawtoken{} deposit, refunded upon approval.

\paragraph{4.\ Service Payments.}
The agent-to-agent economy uses \clawtoken{} as the unit of account for data, compute, and skill services. Representative price ranges include premium skills and tools (10--100~\clawtoken{}), data access (1--50~\clawtoken{}), and compute jobs (5--200~\clawtoken{}).

\subsection{Economic Security}
\label{sec:econ-security}

\subsubsection{Supply Control}

The genesis supply is hard-capped at 1~billion \clawtoken{}. New tokens are created only through the inflation mechanism for validator rewards. The inflation rate converges to a 2\% perpetual floor, providing a sustainable security budget.

\subsubsection{Deflationary Mechanisms}

Future community governance proposals may introduce deflationary mechanisms including:
\begin{itemize}[nosep]
    \item Transaction fee burning (if the gas model evolves from zero-fee)
    \item Service marketplace fee burning
    \item Treasury surplus buyback and burn programs
\end{itemize}

\subsubsection{Whale Risk Mitigation}

Several mechanisms prevent concentration of economic power:
\begin{itemize}[nosep]
    \item No single airdrop allocation exceeds 1\% of total supply (10M cap).
    \item Governance weighting emphasizes reputation and contributions, not only stake (\Cref{sec:governance}).
    \item Quadratic voting for high-impact proposals is planned for future implementation.
\end{itemize}

\subsection{Fair Launch Principles}
\label{sec:fairlaunch}

\claw{} adheres to fair launch principles:
\begin{itemize}[nosep]
    \item \textbf{No pre-mine:} Founding contributors receive 10\%, fully vested over 24~months.
    \item \textbf{No venture capital allocation:} The project is community-driven, prioritizing fairness over capital raises.
    \item \textbf{Full transparency:} All allocations are published on-chain, airdrop snapshots are publicly auditable, and vesting schedules are enforced by smart contracts.
\end{itemize}
